            The {\bf WordSpec} class provides information about the
            mapping of data types into variables. Data types are mapped
            as a 1-dimensional array of words. Each word is a primitive.
            
            A WordSpec may be generated from a Type by the method
            Type.getWordSpec(). There are two signatures, one for the
            entire variable and another for components of the variable
            given subscripts or fields.
            
            The WordSpec class is intended to specify words or bit fields
            for an immediate or target mode variable. These words or bitfields
            should match the components of the type for the variable.

            The member variable {\bf word} is an ArrayList of word indices and is
            used for both immediate mode and target mode variables. When
            the WordSpec applies to an entire 3PL variable this list is a
            contiguous sequence starting at 0. When the WordSpec applies to only a
            part of a variable this list is ascending but not necessarily
            contiguous.
            
            If the WordSpec applies to a target variable there is also a
            matching ArrayList of bit pairs, {\bf bitpairs}, each entry giving the lower
            and upper bit indices of the associated word. For an immediate mode
            variable {\bf bitpairs} is not used. The type of {\bf bitpairs} is
            class {\bf Pair}. If the type is "log" it is not clear if the mode
            is immediate or not so a bitpair 0..0 is included.
            
            Generally a WordSpec describes an entire variable or one or more
            components of a variable. It should not specify a part of a variable
            which does not match the structure of the declared variable.
            The boolean field {\bf complete} is true if a WordSpec specifies a complete
            variable.
            
            A wordSpec can also be used to describe the structure of an expression
            generated by the compiler, which are usually primitive, i.e. have a
            single word and bitpair. These internal signals usually have
            identifiers of the form "E1234" which cannot clash with the
            identifiers of program variables.
            
            Where a portion of a target variable component is required this is
            not obtained by a WordSpec specifying the bits directly as this
            would not match the declaration WordSpec. Instead the entire variable
            component is evaluated and a part is then extracted by a suitable
            operator and connected (assigned) to an internally generated
            signal as described in the previous paragraph.The CONNECT TDE
            can perform some pruning and shifting of bits, otherwise the
            OPERATOR TDE might be used. These methods are used for example
            to extract the sign of a variable or the biased exponent or mantissa
            of a floating point target variable.

            The member variable {\bf types} is an ArrayList of the primitive types
            of the words within the 3PL variable. Each entry in this list
            corresponds to the {\bf word} entry at the same list index.

            Class {\bf Pair} contains four variables -
            \blist
            \item   {\bf lower} - int - the lowest bit offset for this word
            \item   {\bf upper} - int - the highest bit offset for this word
            \item   {\bf indices} - ArrayList - a list of subscripts and field
                    names for this word
            \item   {\bf select} - TDEVar - a select signal used in the case of a
                    target variable subscript
            \blend

            For example the declarations
            \begin{lstlisting}[gobble=12, language=threepl]
                static(s1, "[3][4]int:4");
            \end{lstlisting}
            would create for {\bf s1} a WordSpec containing the lists -

            \begin{tabular}{| l || l || l l | l |}
            \hline
            word  & types    & \multicolumn{2}{c|}{bitpairs} & \\
                  &          & lower & upper & indices \\
            \hline \hline
             0    & typeINT  & 0     & 3     & 0, 0  \\ \hline
             1    & typeINT  & 4     & 7     & 0, 1  \\ \hline
             2    & typeINT  & 8     & 11    & 0, 2  \\ \hline
             3    & typeINT  & 12    & 15    & 0, 3  \\ \hline
             4    & typeINT  & 16    & 19    & 1, 0  \\ \hline
             5    & typeINT  & 20    & 23    & 1, 1  \\ \hline
             6    & typeINT  & 24    & 27    & 1, 2  \\ \hline
             7    & typeINT  & 28    & 31    & 1, 3  \\ \hline
             8    & typeINT  & 32    & 35    & 2, 0  \\ \hline
             9    & typeINT  & 36    & 39    & 2, 1  \\ \hline
            10    & typeINT  & 40    & 43    & 2, 2  \\ \hline
            11    & typeINT  & 44    & 47    & 2, 3  \\ \hline
            \end{tabular}

            which describes the layout of the entire variable. An instance of the
            variable
            \begin{lstlisting}[gobble=12, language=threepl]
             s1[1][1..2]
            \end{lstlisting}
            would create a WordSpec

            \begin{tabular}{| l || l || l l | l |}
            \hline
            word  & types    & \multicolumn{2}{c|}{bitpairs} & \\
                  &          & lower & upper & indices \\
            \hline \hline
             5    & typeINT  & 20    & 23    & 1, 1  \\ \hline
             6    & typeINT  & 24    & 27    & 1, 2  \\ \hline
            \end{tabular}

            which locates two words within the variable.
           
            The entries in {\bf word} must be ascending but do not
            need to be contiguous. An example is a two-dimensional array where
            the second subscript or subscript range does not specifiy the
            entire dimension. If the previous exmaple were changed to -

            \begin{lstlisting}[gobble=12, language=threepl]
             s1[0..1][2..3]
            \end{lstlisting}
            the WordSpec would contain -
            
            \begin{tabular}{| l || l || l l | l |}
            \hline
            word  & types    & \multicolumn{2}{c|}{bitpairs} & \\
                  &          & lower & upper & indices \\
            \hline \hline
             1    & typeINT  & 4     & 7     & 0, 1  \\ \hline
             2    & typeINT  & 8     & 11    & 0, 2  \\ \hline
             5    & typeINT  & 20    & 23    & 1, 1  \\ \hline
             6    & typeINT  & 24    & 27    & 1, 2  \\ \hline
            \end{tabular}

            A further example using both arrays and structs -
            \begin{lstlisting}[gobble=12, language=threepl]
                struct(s,
                    f1, "[4]int:4",
                    f2, "log"
                );
                static(s2, "[3]s");
            \end{lstlisting}
            The WordSpec for the entire variable {\bf s2} would be -

            \begin{tabular}{| l || l || l l | l |}
            \hline
            word  & types    & \multicolumn{2}{c|}{bitpairs} & \\
                  &          & lower & upper & indices \\
            \hline \hline
             0    & typeINT  & 0     & 3     & 0, "f1", 0  \\ \hline
             1    & typeINT  & 4     & 7     & 0, "f1", 1  \\ \hline
             2    & typeINT  & 8     & 11    & 0, "f1", 2  \\ \hline
             3    & typeINT  & 12    & 15    & 0, "f1", 3  \\ \hline
             4    & typeLOG  & 16    & 16    & 0, "f2"     \\ \hline
             5    & typeINT  & 17    & 20    & 1, "f1", 0  \\ \hline
             6    & typeINT  & 21    & 24    & 1, "f1", 1  \\ \hline
             7    & typeINT  & 25    & 28    & 1, "f1", 2  \\ \hline
             8    & typeINT  & 29    & 32    & 1, "f1", 3  \\ \hline
             9    & typeLOG  & 33    & 33    & 1, "f2"     \\ \hline
            10    & typeINT  & 34    & 37    & 2, "f1", 0  \\ \hline
            11    & typeINT  & 38    & 41    & 2, "f1", 1  \\ \hline
            12    & typeINT  & 42    & 45    & 2, "f1", 2  \\ \hline
            13    & typeINT  & 46    & 49    & 2, "f1", 3  \\ \hline
            14    & typeLOG  & 50    & 50    & 2, "f2"     \\ \hline
            \end{tabular}

            An instance
            \begin{lstlisting}[gobble=12, language=threepl]
                s2[1]
            \end{lstlisting}
            would create a WordSpec

            \begin{tabular}{| l || l || l l | l |}
            \hline
            word  & types    & \multicolumn{2}{c|}{bitpairs} & \\
                  &          & lower & upper & indices \\
            \hline \hline
             5    & typeINT  & 17    & 20    & 1, "f1", 0  \\ \hline
             6    & typeINT  & 21    & 24    & 1, "f1", 1  \\ \hline
             7    & typeINT  & 25    & 28    & 1, "f1", 2  \\ \hline
             8    & typeINT  & 29    & 32    & 1, "f1", 3  \\ \hline
             9    & typeLOG  & 33    & 33    & 1, "f2"     \\ \hline
            \end{tabular}

            \begin{lstlisting}[gobble=12, language=threepl]
                s2[1].f1
            \end{lstlisting}
            would create a WordSpec

            \begin{tabular}{| l || l || l l | l |}
            \hline
            word  & types    & \multicolumn{2}{c|}{bitpairs} & \\
                  &          & lower & upper & indices \\
            \hline \hline
             5    & typeINT  & 17    & 20    & 1, "f1", 0  \\ \hline
             6    & typeINT  & 21    & 24    & 1, "f1", 1  \\ \hline
             7    & typeINT  & 25    & 28    & 1, "f1", 2  \\ \hline
             8    & typeINT  & 29    & 32    & 1, "f1", 3  \\ \hline
            \end{tabular}

            and
            \begin{lstlisting}[gobble=12, language=threepl]
                s2[1].f2
            \end{lstlisting}
            would create a WordSpec

            \begin{tabular}{| l || l || l l | l |}
            \hline
            word  & types    & \multicolumn{2}{c|}{bitpairs} & \\
                  &          & lower & upper & indices \\
            \hline \hline
             9    & typeLOG  & 33    & 33    & 1, "f2"     \\ \hline
            \end{tabular}

            Where variable subscripts are used the {\bf select} variable
            in class {\bf Pair} is used. It is only used in a WordSpec
            describing a variable reference or value, not in a Wordspec
            describing a created variable. For example for a variable created
            by -
            \begin{lstlisting}[gobble=12, language=threepl]
             static(s3, "[3][4]int:4");
            \end{lstlisting}
            and evaluated by -
            \begin{lstlisting}[gobble=12, language=threepl]
                static(ss1, "uint:2");
                value(ss2, "uint:2");
                s3[ss1][ss2]
            \end{lstlisting}
            the WordSpec for this evaluation of {\bf s3} might be -

            \begin{tabular}{| l || l || l l | l | l |}
            \hline
            word  & types    & \multicolumn{2}{c|}{bitpairs} & & \\
                  &          & lower & upper & indices & select \\
            \hline \hline
             0    & typeINT  & 0     & 3     & 0, 0 & S147 \\ \hline
             1    & typeINT  & 4     & 7     & 0, 1 & S148 \\ \hline
             2    & typeINT  & 8     & 11    & 0, 2 & S149 \\ \hline
             3    & typeINT  & 12    & 15    & 0, 3 & S150 \\ \hline
             4    & typeINT  & 16    & 19    & 1, 0 & S151 \\ \hline
             5    & typeINT  & 20    & 23    & 1, 1 & S152 \\ \hline
             6    & typeINT  & 24    & 27    & 1, 2 & S153 \\ \hline
             7    & typeINT  & 28    & 31    & 1, 3 & S154 \\ \hline
             8    & typeINT  & 32    & 35    & 2, 0 & S155 \\ \hline
             9    & typeINT  & 36    & 39    & 2, 1 & S156 \\ \hline
            10    & typeINT  & 40    & 43    & 2, 2 & S157 \\ \hline
            11    & typeINT  & 44    & 47    & 2, 3 & S158 \\ \hline
            \end{tabular}

            where the select signals are outputs of decoders for all
            pairs of values of the two variable subscripts.

            If the first subscript was an immediate value -
            \begin{lstlisting}[gobble=12, language=threepl]
                s3[1][ss2]
            \end{lstlisting}
            the WordSpec for this evaluation of {\bf s3} might be -

            \begin{tabular}{| l || l || l l  |l | l |}
            \hline
            word  & types    & \multicolumn{2}{c|}{bitpairs} & & \\
                  &          & lower & upper & indices & select \\
            \hline \hline
             4    & typeINT  & 16    & 19    & 1, 0 & S203 \\ \hline
             5    & typeINT  & 20    & 23    & 1, 1 & S204 \\ \hline
             6    & typeINT  & 24    & 27    & 1, 2 & S205 \\ \hline
             7    & typeINT  & 28    & 31    & 1, 3 & S206 \\ \hline
            \end{tabular}

            where the select signals decode the single variable subscript.

            If the second subscript was an immediate value -
            \begin{lstlisting}[gobble=12, language=threepl]
                s3[ss1][1]
            \end{lstlisting}
            the WordSpec for this evaluation of {\bf s3} might be -

            \begin{tabular}{| l || l || l l | l | l |}
            \hline
            word  & types    & \multicolumn{2}{c|}{bitpairs} & & \\
                  &          & lower & upper & indices & select \\
            \hline \hline
             1    & typeINT  & 4     & 7     & 0, 1 & S341 \\ \hline
             5    & typeINT  & 20    & 23    & 1, 1 & S342 \\ \hline
             9    & typeINT  & 36    & 39    & 2, 1 & S343 \\ \hline
            \end{tabular}

            The select signals are mutually exclusive, each being high when the
            value of the subscript variables matches the indices of the associated
            word (array member). For a variable evaluation these signals are used
            with a SELECT TDE element to multiplex the output of a variable. In
            the case of a variable assignment using variable subscripts these
            signals are used to select the word within the variable which is to be
            written.

            The following methods for class WordSpec provide access
            to information for the member variables discussed above -
            \blist
            \item   numWords() - the number of words
            \item   numBits() - the total number of bits
            \item   getWord(i) - the ith word index
            \item   getPrimType(i) - the ith word primitive type
            \item   getLower(i) - the ith word lower bit index
            \item   getUpper(i) - the ith word upper bit index
            \item   getWidth(i) - the ith word width (bits)
            \item   getSelect(i) - the ith word select signal
            \blend

            Other member variables for WordSpec are -

            {\bf ArrayList subvals}

            If this word specification applies to an expression containing
            target variable subscripts then this list contains the {\bf Val}s for
            those subscripts in left-to-right order. The {\bf select} entries
            in the {\bf Pair} list are filled in by a single call to method
            {\bf processVarSubs()} which uses the entries in member variable
            {\bf subvals}.

            {\bf int[] dim\_des}

            A dimensional description - see below.

            {\bf Type check\_type}

            The type to check for assignment - see below.

            {\bf Type type}

            The type - see below.

            {\bf QueueRefs squeues}

            If this word specification applies to an expression containing
            queue reads or to a queue variable being written then this member variable
            specifies the queue dependencies.

            {\bf Var var}

            If this word specification applies to a variable then this class member
            is the variable.



            \subsubsection{the dim\_des, check\_type and type fields}

                When a variable is assigned a value the type of the variable
                must be checked against the type of the value. For primitive
                types this is straightforward. For array types the process
                is more complicated. For example the 3PL code
                \begin{lstlisting}[gobble=16, language=threepl]
                    var(a, "[7][3][4]int");
                    var(b, "[3][4]int");
                    a[2..4][1] = b;
                \end{lstlisting}
                is correct in that the dimensionality of {\bf a[2..4][1]} is
                {\bf [3][4]}. This is because the first subscript is a subrange
                of dimension 3, the second subscript eliminates that dimension
                and the third subscript is missing and so defaults to the
                dimension {\bf [4]}.

                When this class was first implemented the view was taken that
                the problem of dimensionality could best be solved by dividing
                a reference into leading subscripts followed by a type. The
                only complicating issue is the occurence of structs nested
                among arrays. To access an array within a struct however requires
                that the struct field be explicitly given, which immediately
                eliminates the rest of the struct. On the other hand any struct
                which is referenced intact cannot have the dimensionality of
                nested arrays changed (subrange or single subscript). Thus all
                references no matter how complicated can be reduced for the
                purposes of dimensionality and type checking into a leading
                number of array dimensions followed by a type. The type
                may itself contain arrays, however the dimensions of these
                will be as created, i.e. have not been changed by subscripts
                or subscript ranges. The dimensional description is the field
                {\bf dim\_des} which is a one-dimensional array whose size
                gives the number of dimensions of the reference and each entry
                gives the dimension. The type is the {\bf check\_type} field
                which can be checked for equality. The dimensional description
                describes the reference through the leading subscripts, be they
                missing dimensions or sub arrays including struct fields of that
                type. Single subscripts do not appear as they eliminate that
                dimension, similarly struct field references disappear as they
                eliminate the struct as a whole. When the description finally
                encounters a primitive of a complete struct (regardless of what
                the struct contains, including nested arrays) the rest of
                the type becomes {\bf check\_type}. For examples see appendix A.

                It was later found necessary to add a proper type field representing
                the actual type of the reference. This is easily generated by
                iterating through the {\bf dim\_des} array creating nested
                class {\bf Type} of type array and then finally nesting
                {\bf check\_type} at the end.

                In retrospect it would have been easier and simpler to generate
                the type field directly and check types against each other for
                equality rather than bother with the dimensional description and
                then have to generate {\bf Type} from it anyway. The code could be
                changed, however the tentacles of the current approach spread
                widely into surrounding code and eliminating bugs introduced by
                such a change would take some time. It may be cleaned up in the
                future but for the moment this obscure mechanism remains in place.
